% ====================================================================
% DATEI: content/08_boundary.tex
% Erweiterte Rand- und Patch-Geometrie — Core 1.3 (Abschlussversion)
% ====================================================================

\section{Erweiterte Rand- und Patch-Geometrie}
\label{sec:boundary-extended}

Dieses Kapitel stellt die vollständige Theorie von Grenzen und Patch-Geometrie wieder her.
Ränder bilden eine der universellsten Strukturen im OC-Rahmen:
Sie steuern die Phasentrennung in \(K_2\), die Membranbildung in \(K_3\)--\(K_4\),
die Erregbarkeit in \(K_5\) sowie repräsentationale oder institutionelle Grenzen in
\(K_6\)--\(K_8\).
Die vorliegende Version konsolidiert die für Core~1.3 erforderliche vollständige Struktur:
schwellenwertdefinierte Grenzen, Patch-Diskretisierung, lokale Schwellenwertklassen
und randgetriebene Übergänge.

\subsection{Formale Definition von \texorpdfstring{$\partial\Omega(K)$}{boundary}}

Für jedes Kontinuum \(K\) mit zulässigem Zustandsraum \(\Omega(K)\) und Schwellenwertlandschaft
\(\Theta(K)\) ist der Rand als die Menge von Zuständen definiert, für die
mindestens ein Schwellenwert exakt gesättigt ist:
\[
  \partial\Omega(K)
  =
  \big\{\, s \in \overline{\Omega(K)} :
          \exists k\ \text{such that}\ g_k(s) = \Theta_k \big\}.
\]
Äquivalent dazu gilt: Wenn \(f_k(s)=g_k(s)-\Theta_k\) die verschobene Restfunktion
bezeichnet, die im Notizanhang der Notation und im Modellkapitel verwendet wird,
halten Innenzustände \(f_k(s)<0\), Zustände auf der zulässigen Seite \(f_k(s)\le0\), und
der gleiche Rand ist die Nullstelle \(f_k(s)=0\) für mindestens einen aktiven Schwellenwert.

Damit gilt:
\begin{itemize}
    \item Innenzustände erfüllen \(f_k(s) < \Theta_k\) für alle Schwellenwerte;
    \item Außenzustände verletzen mindestens einen Schwellenwert;
    \item der Rand ist eine \emph{schwellenwertdefinierte Oberfläche}, keine geometrische
          Primitive.
\end{itemize}

Diese Definition vereint klassische Oberflächen, Membranen, Vesikel,
Perkolationsfronten, erregbare Grenzen, logische Limiten und institutionelle
Schwellenwerte in einem Rahmen.

\subsection{Randgeometrie}

Randgeometrie ergibt sich aus der lokalen Struktur von Schwellenwerten, Potentialen
und Flüssen.
Wenn der aktive Schwellenwert im Punkt \(s \in \partial\Omega(K)\) \(g_k\) ist, ist die
äußere Normale parallel zur Gradientenrichtung:
\[
  n(s) \propto \nabla g_k(s).
\]

Mehrere aktive Schwellenwerte erzeugen Ecken, Kanten und Regionen hoher Krümmung.
Krümmung wird dynamisch relevant, sobald Biege- oder Oberflächenenergie-Terme
in das Potential eingehen:
\[
  P_{\mathrm{bend}}(s) = \tfrac{1}{2}\,\kappa(s)^2\,\kappa_b,
\]
mit Biegesteifigkeit \(\kappa_b\).

Dies verbindet:
\begin{itemize}
    \item \(K_2\): klassische Oberflächenspannung,
    \item \(K_3\)–\(K_4\): Membran-Krümmung und Vesikel-Energetik,
    \item \(K_5\): Geometrie erregbarer Grenzen.
\end{itemize}

\subsection{Patch-Modell von Grenzen}

OC modelliert Grenzen nicht als kontinuierliche, homogene Flächen, sondern als
Sammlungen dynamischer \emph{Patches}:
\[
  \partial\Omega(K) = \bigcup_i \partial\Omega_i.
\]

Jeder Patch \(P_i\) enthält:
\[
  (\sigma_i,\Theta_i,P_i(t),J_i(t)),
\]
was seinen lokalen Zustand, die Schwellenwerte, Potentiale und Flüsse beschreibt.

\paragraph{Lokale Zustände.}
Patches nehmen symbolische Zustände an:
\[
  \sigma_i \in \{L_\alpha, L_\beta, L_o, L_f, L_b\},
\]
wobei sie lokale Packungs- oder Strukturregime kodieren:
\begin{itemize}
    \item \(L_\alpha\): flüssig-ungeordnet,
    \item \(L_\beta\): geordnet/gelartig,
    \item \(L_o\): flüssig-geordnet,
    \item \(L_f\): fragil/dünn (prärupturisch),
    \item \(L_b\): gebrochener Randzustand.
\end{itemize}

Dieses Alphabet ist über biologische Membranen hinaus verallgemeinerbar auf
Perkolationskanten, Rissfronten, fluktuierende Interfaces und sogar symbolische/kognitive
Grenzen.

\paragraph{Lokale Schwellenwerte.}
Jeder Patch besitzt einen Vektor:
\[
  \Theta_i = (\Theta^{\rm perm}_i, \Theta^{\rm grad}_i, \Theta^{\rm mem}_i,
              \Theta^{\rm bend}_i, \ldots ),
\]
mit den zugehörigen lokalen Potentialen \(P_i(t)\) und Flüssen \(J_i(t)\).

\subsection{Grenzwertsystem}

Über \(K_2\)--\(K_5\) erscheinen universell drei Schwellenwertklassen.

\paragraph{1. Permeabilitätsschwelle \(\Theta_{\rm perm}\).}
Ein Patch wird permeabel, wenn:
\[
  |\nabla P_i| > \Theta^{\rm perm}_i.
\]
Beispiele:
\begin{itemize}
    \item osmotische Permeabilität (\(K_3\)--\(K_4\)),
    \item Ionenleck-Beginn (\(K_5\)),
    \item Phasenfrontdurchdringung (\(K_2\)).
\end{itemize}

\paragraph{2. Gradienten-Schwelle \(\Theta_{\rm grad}\).}
Wenn ein Gradient nicht aufrechtzuerhalten ist:
\[
  |\nabla P_i| > \Theta^{\rm grad}_i
  \quad\Rightarrow\quad \sigma_i \to L_f.
\]
Dies steuert:
\begin{itemize}
    \item Vesikelriss bei hohem osmotischem Druck (\(K_4\)),
    \item Auslösung von Aktionspotenzialen an Initialsegmenten von Axonen (\(K_5\)),
    \item beschleunigte Phasenfronten und Rissausbreitung (\(K_2\)).
\end{itemize}

\paragraph{3. Membranintegritätsschwelle \(\Theta_{\rm mem}\).}
Wenn Membranspannung oder Krümmung Grenzwerte überschreiten:
\[
  \gamma_i > \Theta^{\rm mem}_i
  \qquad\text{oder}\qquad
  \kappa(s) > \Theta^{\rm bend}_i,
\]
dann gilt:
\[
  \sigma_i \to L_b.
\]

Dies liefert ein vereinheitlichtes Verständnis von Rissen, Beulenbildung, Porenbildung,
Randversagen und Oberflächenkollaps.

\subsection{Randgetriebene Dynamik}

Nachbarschaftliche Patches interagieren über Flüsse:
\[
  J_{ij} = F(P_i, P_j, \sigma_i, \sigma_j).
\]

Damit entstehen:
\begin{itemize}
    \item Stabilisierung und Umordnung von Patch-Mosaiken,
    \item Ausbreitung von Rissfronten,
    \item krümmungsgetriebene Migration,
    \item Porenbildungseinleitung,
    \item Verhalten als erregbares Medium in \(K_5\).
\end{itemize}

Patch-Entwicklung folgt:
\[
  \frac{d\sigma_i}{dt}
  = S(\sigma_i, P_i, J_i, \Theta_i),
\]
wohin \(S\) ein struktureller Aktualisierungsoperator ist, unabhängig vom
betroffenen Definitionsbereich.

\medskip
\textbf{Prinzip der Randlokalisation.}
\emph{Alle Dimensional-Übergänge in \(K_3\)--\(K_5\) beginnen an Randpatches.}

Dies erklärt:
\begin{itemize}
    \item den Protocell-Rupturbeginn an Hoch-Krümmungsstellen,
    \item die Auslösung von Aktionspotenzialen an den Axoninitialsegmenten,
    \item Phasenübergänge, die an Defekten oder Oberflächen nucleieren.
\end{itemize}

\subsection{Randversagen und Kollaps}

Randkollaps tritt ein, wenn eine zusammenhängende Menge von Patches gebrochen wird:
\[
  \sigma_i = L_b
  \qquad \forall\, i \in \mathcal{C},
\]
und der Cluster \(\mathcal{C}\) sich über den Rand perkoliert.

Nach Satz~3 (Ergebniskapitel):
\[
  \partial\Omega(K)\ \text{kollabiert}
  \quad\Longrightarrow\quad
  \Omega(K) = \emptyset,
\]
und das Kontinuum erlischt.

Wir unterscheiden:
\begin{itemize}
    \item \textbf{lokales Versagen}: isolierte \(L_b\)-Patches, reversibel;
    \item \textbf{Cluster-Kollaps}: perkolierendes \(L_b\)-Netzwerk, irreversibel;
    \item \textbf{globaler Randverlust}: Verschwinden des Randes, was Tod oder Übergang
          zu einem neuen Kontinuum impliziert.
\end{itemize}

\subsection{Rand und Wiedergeburt}

Die Entstehung eines neuen Kontinuums \(K_{x+1}\) erfordert typischerweise:
\begin{itemize}
    \item Bildung oder Riss von Randpatches,
    \item Erzeugung neuer Patch-Typen,
    \item Auftreten neuer Schwellenwertklassen an Randbereichen,
    \item Etablierung einer neuen globalen Randgeometrie.
\end{itemize}

\textbf{Prinzip der Randwiedergeburt.}
\emph{Neue Kontinua entstehen aus neuen Randgeometrien.}

Beispiele:
\begin{itemize}
    \item Membranschluss im Übergang \(K_3 \to K_4\),
    \item Auftreten erregbarer Grenzen im Übergang \(K_4 \to K_5\),
    \item Bildung institutioneller Grenzen im Übergang \(K_6 \to K_7\),
    \item Auftreten repräsentationaler Grenzen im Übergang \(K_9 \to K_{10}\).
\end{itemize}

Grenzen sind daher nicht nur Beschränkungen:
Sie sind generative Oberflächen, auf denen neue Achsen, Restriktionen und Zyklen entstehen.
